\scanpage{19}
% archive foliation: 18
\begin{center}- 90 -\end{center}\hfill n\textsuperscript{o} 268

nous pouvons supposer que $P\notin\operatorname{supp}(D)$. Soit $f_0,\ldots,f_n$ une base de $L(D)$ telle que toutes les fonctions $f_i$ soient définies sur $K$. Toute fonction $f\in L(D)$ peut être écrite
\[
f=f_\lambda=\sum_{i=0}^n\lambda_i f_i\qquad\lambda_i\in\Omega\,.
\]
Puisque aucun des $f_i$ n'a de pôle passant par $P$, chaque $f_i$ est dans $\unsure{\mathfrak{o}|v}$ (IAG) et a une expension en série de puissance
\[
f_i=\sum_{\mu=0}^{\infty}c_{i\mu}M_\mu(t)\qquad c_{i\mu}\in K
\]
et par conséquent
\[
f_\lambda=\sum_{\mu=0}^{\infty}\left(\sum_{i=1}^n\lambda_i c_{i\mu}\right)M_\mu(t).
\]
La matrice infinie $\|c_{i\mu}\|$ est de rang $n+1$ parce que les $f_i$ sont linéairement indépendantes sur $K$. Soit $s$ le plus petit entier tel que $\|c_{i\mu}\|\quad i=0,\ldots,n\quad\text{et}\quad\mu=0,\ldots,s$ soit de rang $n+1$. Alors les équations
\[
\sum_{i=0}^n\lambda_i c_{i\mu}=0\qquad\mu=0,\ldots,s
\]
déterminent les $\lambda_i$ de façon unique dans $K$ (à un facteur d'homogénéité près). La fonction correspondante $f_\lambda$ a un diviseur
\[
(f_\lambda)=X-D\qquad X>0.
\]
Puisque $f_\lambda$ est rationnelle sur $K$, il s'ensuit que $X$ est rationnel sur $K$. Il est clair d'après la construction que le diviseur $X$ est le diviseur dont on a démontré l'unicité dans (i). Ceci démontre notre \deleted{théorème.} \inserted{\unsure{prop}}

\uline{Théorème} - Soit $D$ un diviseur positif sur un produit $A\times V$ d'une variété de groupe commutatif et d'une variété $V$ \uline{complète sans singularités de codimension 1}. Soit $k$ un corps de définition pour $A$ et $V$, sur lequel $D$ soit rationnel, et tel que $V$ possède \uline{un point simple $P$ rationnel sur $k$}. Alors il existe une variété de groupe commutatif $B$ et un
\margin{\unsure{utile au cas non singulier?}}
\margin{\unsure{Nouveau!|Revoir!}}
